% ============================================================
%  CV -- Loujein Bousnina
%  Cible : poste en systemes embarques / IoT / electronique
%  Style : CV compact, serif classique (Latin Modern)
%  Compilation : pdflatex main.tex
% ============================================================
\documentclass[10pt,a4paper]{article}

\usepackage[a4paper,top=1.3cm,bottom=1.3cm,left=1.8cm,right=1.8cm]{geometry}
\usepackage[T1]{fontenc}
\usepackage[utf8]{inputenc}
\usepackage{lmodern}
\usepackage[french]{babel}
\usepackage{microtype}
\usepackage{tabularx}
\usepackage{enumitem}
\usepackage{titlesec}
\usepackage{xcolor}

\definecolor{linkblue}{HTML}{1155CC}

\usepackage{hyperref}
\hypersetup{
  colorlinks=true,
  urlcolor=linkblue,
  linkcolor=linkblue,
  pdftitle={CV - Loujein Bousnina - Systèmes embarqués et IoT},
  pdfauthor={Loujein Bousnina}
}

\pagestyle{empty}
\setlength{\parindent}{0pt}

% ---------- Titres de section ----------
\titleformat{\section}{\large\bfseries\scshape}{}{0em}{}[\vspace{1pt}\titlerule]
\titlespacing*{\section}{0pt}{8pt}{4pt}

% ---------- Listes ----------
\setlist[itemize]{leftmargin=12pt, itemsep=0pt, topsep=1pt, parsep=0pt, label=\textendash}

% ---------- Ligne Formation : date | Ecole --- Diplome | note (droite, italique) ----------
\newcommand{\cvrow}[4]{%
  \noindent\begin{tabularx}{\linewidth}{@{}p{2.3cm}X@{}}
    \small #1 & \textbf{#2} --- #3 \hfill\textit{\small #4}\\
  \end{tabularx}\par\vspace{3pt}%
}

% ---------- Ligne Experience / Projet : Titre (gauche) -- Date (droite) ----------
\newcommand{\cvhead}[2]{%
  \noindent\begin{tabularx}{\linewidth}{@{}X r@{}}
    \textbf{#1} & \textit{\small #2}\\
  \end{tabularx}\vspace{1pt}%
}

\begin{document}

% ================= EN-TETE =================
\begin{center}
  {\LARGE Loujein Bousnina}\\[3pt]
  {\large Systèmes embarqués \textbullet{} IoT \textbullet{} Électronique}\\[4pt]
  {\small +216 24 411 846
    \quad$|$\quad \href{mailto:loujeinbousnina2020@gmail.com}{loujeinbousnina2020@gmail.com}
    \quad$|$\quad \href{https://www.linkedin.com/in/loujein-bousnina/}{linkedin.com/in/loujein-bousnina}}\\[1pt]
  {\small Bardo, Tunisie \quad$|$\quad Né le 04/11/2000 \quad$|$\quad Permis B}
\end{center}
\vspace{3pt}
\hrule
\vspace{2pt}

% ================= PROFIL =================
\section{Profil}
{\small
Jeune diplômée du cycle d'ingénieur en informatique, spécialisée en Systèmes Logiciels Embarqués, Ambiants
et Mobiles (SLEAM). Passionnée par les nouvelles technologies, l'IoT, l'intelligence artificielle et les systèmes
embarqués. Recherche un poste en systèmes embarqués, IoT ou électronique.
\par}

% ================= FORMATION =================
\section{Formation}
\cvrow{2023 -- 2026}{Esprit, Ariana}{Cycle d'ingénieur en informatique, spécialité \textbf{SLEAM}\newline
  \textit{\small Systèmes Logiciels Embarqués, Ambiants et Mobiles}}{}
\cvrow{2019 -- 2023}{ISTIC, Borj Cédria}{Licence en Informatique Industrielle}{}
\cvrow{2019}{Lycée Sijoumi, Tunisie}{Baccalauréat Sciences Techniques}{}

% ================= COMPETENCES =================
\section{Compétences techniques}
\noindent\begin{tabularx}{\linewidth}{@{}>{\bfseries}p{5.2cm}X@{}}
  Embarqué \& électronique    & \small ESP32, Arduino, Raspberry Pi, C/C++ embarqué \\[2pt]
  IoT \& communication        & \small MQTT, GSM, RFID/NFC, GPS, Node-RED \\[2pt]
  Langages \& IA              & \small C, C++, Python, JavaScript, TypeScript ; modèles LSTM \\[2pt]
  Web, mobile \& DevOps       & \small React, React Native, Flutter, Node.js, MongoDB, Docker \\[2pt]
  Savoir-être                 & \small Autonomie, organisation, communication, leadership \\
\end{tabularx}

% ================= EXPERIENCE =================
\section{Expérience professionnelle}

\cvhead{MakerSkills --- Stage PFE : Développeur IoT \& Full-Stack}{Fév. 2026 -- Août 2026}
\begin{itemize}
  \item \textbf{RondeGuard} : solution IoT de suivi d'agents en temps réel (\textbf{ESP32, MQTT, RFID/NFC}).
  \item Applications mobile et web (React Native, React, Node.js, MongoDB, Docker).
\end{itemize}

\cvhead{SNCFT, Département Informatique --- Stage ingénieur}{Juin 2025 -- Sept. 2025}
\begin{itemize}
  \item Système de \textbf{localisation des trains en temps réel} par GPS et IoT.
  \item ESP32, Node-RED et MongoDB pour le traitement et la visualisation des données.
\end{itemize}

\cvhead{SNCFT, Département Informatique --- Stage}{Juin 2023 -- Août 2023}
\begin{itemize}
  \item Optimisation \textbf{SEO} et indexation des pages web pour améliorer la visibilité en ligne.
\end{itemize}

\cvhead{SOTRADIES, Département Technique --- Stage PFE licence}{Fév. 2023 -- Mai 2023}
\begin{itemize}
  \item Remise en état de la commande d'un pont élévateur à 4 colonnes.
\end{itemize}

\cvhead{Chakira Cable, Maintenance générale --- Stage technicien}{Juil. 2022 -- Août 2022}
\begin{itemize}
  \item Conception d'un système de sécurité pour armoire électrique à base de microcontrôleurs \textbf{Arduino}.
\end{itemize}

\cvhead{Coficab, Département Maintenance --- Stage}{Juil. 2021 -- Août 2021}
\begin{itemize}
  \item Maintenance d'équipements électriques industriels.
\end{itemize}
% ================= PROJETS =================
\section{Projets académiques}

\cvhead{WaterSec (Istidema) --- Projet intégré, Esprit}{Fév. 2025 -- Juin 2025}
\begin{itemize}
  \item Solution intelligente de gestion de l'eau (application Flutter, MongoDB).
  \item Détection de fuites par \textbf{LSTM} (85\,\% de précision) ; \textbf{$-$30\,\% de gaspillage d'eau}.
\end{itemize}

\cvhead{Smart Health Glove --- Gant connecté de santé (wearable)}{}
\begin{itemize}
  \item Gant intelligent de suivi de la santé, avec antenne intégrée pour la communication sans fil.
\end{itemize}

% Format pour ajouter un projet :
% \cvhead{Titre du projet --- type}{2025}
% \begin{itemize}
%   \item Objectif du projet et votre rôle.
%   \item Matériel et technologies : carte, capteurs, protocoles (I2C, SPI, UART...).
%   \item Résultat concret (mesure, démonstration, note...).
% \end{itemize}

% ================= CERTIFICATIONS =================
\section{Certifications}
\noindent\begin{tabularx}{\linewidth}{@{}>{\bfseries}p{5.2cm}X>{\itshape\small}r@{}}
  AWS Academy           & \small Cloud Foundations & oct. 2025 \\[2pt]
  Hashgraph Association & \small Hashgraph Developer (attestation de participation) & oct. 2025 \\[2pt]
  NVIDIA                & \small AI on Jetson Nano \textbullet{} AI for Anomaly Detection \textbullet{} AI for Predictive Maintenance \textbullet{} Generative AI with Diffusion Models & janv. 2025 \\[2pt]
  DataCamp              & \small AI Fundamentals \textbullet{} AI Ethics & janv. 2025 \\
\end{tabularx}

% ================= LANGUES =================
\section{Langues}
\small Arabe (langue maternelle) \quad$|$\quad Français (courant) \quad$|$\quad Anglais (courant) \quad$|$\quad Allemand (débutant)

\end{document}
